% !TeX root = ../draft_main.tex
\section{Chaometer as a measure of "directionality"}

For the particular case of the bipartition of one spin and the rest of the chain in a spin model, which we well be calling the chaometer for short, the relationship between the purity $P(t)$, entanglement (measured by the linear entropy) $S_{L}(t)$ and the direction in the Bloch sphere of the chaometer is one to one

\begin{equation}
    P(t) = Tr_{L-1}\{\hat{\rho}_{1}(t)^2 \} = 1 - S_{L}(t) = \dfrac{1+\langle\sigma_{x}\rangle+\langle\sigma_{y}\rangle+\langle\sigma_{z}\rangle}{2} .
    \label{eq:FUNDAMENTAL}
\end{equation}

With the condition that we are always starting from a pure state $\ket{\psi(0)} = \ket{\psi_{1}}\otimes\ket{\psi_{L-1}}$ and evolving it with an unitary operator $\hat{\rho}_{1}(t) = Tr_{L-1}\{\hat{U}(t)\ket{\psi(0)}\bra{\psi(0)}\hat{U}^{\dagger}\}$. There are several remarks from the previous expression. \underline{Number one} , the purity can not reach the minimum value 1/2, because the initial state we are considering is always a product state between the bipartition, the only way the purity achieves the maximum value is in the case that the initial state is a Bell state\janote{The state after the evolution can become entangled}. However, the value can get very close to 1/2 (CITE OR SHOW) with the dimension of the chain, as the Page value, for the von Neumann entropy, sets a value for the average entanglement obtained from a random state in the Full Hilbert space (not to confuse with the maximally entangled state). \underline{Number two}, with only the value of the purity one can not say anything about the overall direction of the Bloch vector of the chaometer\janote{I agree}. \underline{Number three}, the initial decay of the purity can be analytically obtained (CITE PAPER, universal decay purity something...) and it goes like $t^{-2}$ (exactly the same reason as the survival probability)\janote{Ok. so you're saying it depends only on the density of states}. Other than those previous remarks, the purity can stay constant, decay and increase as a function of time without any restriction due to the fact that we are not limiting to any type of Markovian dynamics, which would imply a monotonous decay to reach a steady state\janote{However, due to ergodicity one must have evolved states that are approximately Haard random after some time}. \underline{Number four}, the previous remarks were made having in mind a single initial state. If one is observing the average purity of an ensemble of states, then the remarks need to be taken as properties of the ensemble (at this moment, for me, taking an average is not mathematically the same as being independent, in the sense of the super-operator being independent of the state of the chaometer)\janote{I didn't understand this point}.

\textbf{Hypothesis}: If there is not a privileged direction from the Hamiltonian itself, $\hat{U}(t) = exp[-i\hat{H}t]$, then the purity of the chaometer will decay to its lowest possible value with almost none fluctuations as a function of time.\\

\emph{Evidence}: By only observing the purity of the chaometer, one does not have access to the overall direction of the dynamics of the Bloch vector of the chaometer. If the purity shows persistent fluctuations over time, by \ref{eq:FUNDAMENTAL}, one or more of the components of the Bloch vector are necessarily fluctuating. To the contrary, if the fluctuations are very small (close to 1/2 by remark one), then the Bloch vector is necessarily confined to a tiny neighborhood around the origin, very close to the effect of a completely depolarizing channel, which has no preferred direction. \janote{This 
sounds that can be argued from the previous ergodic argument}

In my opinion, all the observations that we have made, our calculations (endless amount of calculations, which their usefulness are only visible now, even though I will not include the vast majority here) and results from other papers, are explained by the hypothesis and the evidence. A priori, there is no need to introduce chaos in the observed behavior of the purity of the chaometer. If the particular Hamiltonian in question happens to have quantum chaotic properties, then, of course the purity will be linked to the chaotic properties of said Hamiltonian, but as I will explain, the best case for the purity as a quantum chaos indicator is a proxy for quantum chaos, very similarly to the IPR\janote{proxy?}\footnote{The IPR can be fooled too. Jonathan showed in the last year in Xalapa that, the IPR of the eigenvectors of a spin chain with \underline{correlated} noise in the computational basis, can mistakenly show localization while spectral statistics show quantum chaos.} (this last statement can be redefined in order to fit other narrative). For me, it can not be considered at the same level of the $\langle r\rangle$ or survival probability (this last statement can be redefined in order to fit a more optimistic narrative).\janote{Maybe we can say something about the long-range correlations in the choi purity}\\

Let me explain what do I mean by "privileged direction from the Hamiltonian". For the chaometer it is very clear, the Bloch vector of the density matrix defines what I mean by "direction". For the Hamiltonian, let me write the most general Hamiltonian that we have read or used (with open boundary conditions for simplicity now), which only contains nearest neighbor interactions.

\begin{equation}
\begin{split}
\hat{H} &= \sum_{i=1}^{L-1}(a\hat{\sigma}_{i+1}^{x}\hat{\sigma}_{i+1}^{x} + b\hat{\sigma}_{i+1}^{y}\hat{\sigma}_{i+1}^{y} + c\hat{\sigma}_{i+1}^{z}\hat{\sigma}_{i+1}^{z}) \\
  &\quad + \sum_{i=1}^{L}(d\hat{\sigma}_{i}^{x} + e\hat{\sigma}_{i}^{y} + f\hat{\sigma}_{i}^{z})+\epsilon_d\hat{\sigma}_\alpha .
    \label{eq:generalhamiltonian}
\end{split}
\end{equation}

The notion of privileged direction from the Hamiltonian comes from the numerical value of the coefficients. If all of the coefficients are equal or roughly equal in magnitude\footnote{I am not claiming that this is the most isotropic Hamiltonian, this expression is just for you to see the impact of the coefficients compared to the defect. Also, the most isotropic Hamiltonian depends on the model, which is exactly the most chaotic point in parameter space. For the case of the XXZ integrable, the most isotropic case is the one which gives a decaying purity like the chaotic ones.}, then I will call that Hamiltonian isotropic. If there is a single coefficient which is bigger than the others, then I will call that Hamiltonian as having a preferred direction, not necessarily the same direction from the particular term of the Hamiltonian.
The most notorious term in the last expression \ref{eq:generalhamiltonian} is the last one, which represents a magnetic field in a single spin of the whole chain, or a defect as Lea calls it. This last term is completely different from the rest, because if the last coefficient $\epsilon_d$ is higher or lower than the other six parameters, then that does not remove the isotropy of the Hamiltonian immediately (as seen by the chaometer), this is because it is only acting in a single spin. Meanwhile, notice the other six parameters, they are acting on the whole chain, even in the case of being taken from a random distribution as is the case of the models with random noise. So, the overall effect of the defect can be tricky, and this distinction alone comes from the single model we have analyzed.\\

In the previous case, if the model happens to be chaotic, it has always been the case that, the chaotic regime always occurs when the parameters are very close to each other, in other words, when the Hamiltonian is isotropic, so it has no privileged direction, at least in all the models that we have used. When the Hamiltonian transitions into a regular regime, we numerically do this by increasing or decreasing one of the parameters. Think about that for a moment. So, even though the particular physics of the regular regimes are completely different and are system dependent (free fermions for Ising, many-body/Anderson localization for Heinsenberg with noise), from the "directionality point of view" all of the regular regimes are in fact characterized by the dominance of some part of the Hamiltonian over the other, and the only way to achieve that (with these local models) is by lowering or increasing some of the parameter with respect to the other ones.\\

By the same token, the completely integrable by Bethe anzats, the XXZ chain can achieve the same effect as a Hamiltonian with a quantum chaos transition. Needless to say, not starting from the general Hamiltonian \ref{eq:generalhamiltonian}. If you start with the full initial integrable model, it shows in the purity no directionality, so the Hamiltonian in that form is already isotropic. If you change one of its parameters then you create a preferred direction with respect to the original isotropic case, making the purity fluctuate and "mimic" the chaos to regularity transition. This model was one of the key for the hypothesis, because there is no quantum chaos to attribute the behavior of the purity here.\\

What happens for the setup use by Gorin, Pineda, Zirac et al (CITE THE 3-4 PAPERS), which consist on evolving a product state with a unitary from an ensemble of RMT. The global evolution is still unitary, $\hat{U}(t) = exp[-i\hat{H}t]$, and $\hat{H}$ was taken from a GOE ensemble and a Poisson diagonal matrix. In this case, the $\hat{H}$ is completely isotropic, a matrix from RMT ensembles do not posses any other property but the symmetry, or lack of, under time reversal. So, any Hamiltonian from those ensembles, first, can not be thought as a nearest neighbor spin chain or any chain at all. And secondly, if one uses a transition model to go from Wigner-Dyson to Poisson, there is not any direction imprinted in the transition (as confirmed by their results). As concluded here \cite{Gorinpurity2002}, there isn't a distinction from a chaotic or integrable environment. Again, their setup is not exactly the same as ours. It could be debatable if these results do apply for our setup, and my argument is that in this case, the unitary evolution coming from a featureless Hamiltonian, leaves no direction in which a regular regime could appear and make the purity to fluctuate in a greater manner compared to the chaotic case.\\

The last part left to explain is the defect Hamiltonian. In this case, my hypothesis stills holds, but the actual range of values, for which the overall direction the chaometer does not respond is greater compared to the other chaotic models. The reason is very simple, even though the defect is enough to break integrability and induce chaos, in order to produce a privileged direction for the chaometer to detect, the magnitude is larger compared to the other chaotic models. Or at a fixed value for the defect, the other terms needs also acquire larger values to produce a detectable direction for the chaometer to point towards.\\

All the other quantum chaotic Hamiltonians, are easily explained and the parameter range for which the transition from chaotic to regular regime happens (remember in this case, the six parameters affect the whole chain) is also enough to produce a directionality. But, even though, this is the easiest case to explain, there are still a lot of caveats. I have several results for the Ising model which I have studied in more detail than all of the other models, and the purity does not seem to follow a one to one correspondence to the transition marked by $\langle r\rangle$, so, even in the best case scenario which was supposed to be the Ising model for the chaometer.\\

So, the purity of the chaometer (with this narrative) is a proxy for quantum chaos, it responds to that, but the total decoherence, which is the integral of the purity from 0 up to a certain time, can be made to follow the most isotropic case, which I consider as a benchmark the RMT results. In other words, I think that quantum chaotic Hamiltonians, only in the fully chaotic regime, are all universal (this is by no means any new statement) and comparable to RMT ensembles. In this regime, a \underline{generic} product state,as in our case, a product state between every possible bipartition or as Wisniacki says, a spin state at infinite temperature (by tracing out the environment, the reduced density matrix can not achieve the maximally mixed state, simply because the initial state can never be the maximally entangled, why did Wisniacki took the infinite temperature density matrix then?), has contribution among all the eigenvectors making it become very quickly a random vector, making any bipartition to entangle very rapidly to the highest possible value without further changes, saturation/equilibration. Because the global state is almost identical to a random vector (this is why the saturation is always the Page value). If you move in parameter space just slightly, you move away from this \textbf{generic, chaotic and isotropic} regime and a preferred direction is created ultimately making the purity to fluctuate.\\

The final narrative that we can give to these results are completely up to us. I have come up with at least 3 which we can discuss in the next meeting. I also hope to highlight the importance of the whole entanglement spectrum, the Schmidt numbers, or the Bloch components, in order to further distinguish the fluctuations of the purity coming from the three directions, which indicates a kind of isotropy or if it's just one direction. For the Choi state, the 4 eigenvalues will be important too.\\

Is there a way to measure this "directionality" (if there it exist) independently from entanglement properties? Wouldn't it be better to define isotropy directly from the purity in the most chaotic case, put it as a benchmark and then measure deviations from this regime? The most general evolution for a density matrix is a CPTP map, and I think that the most isotropic evolution would be the depolarizing channel, so the directionality can be "tested" as the given channel expressed in the Pauli basis and checked if this is enough to make the purity fluctuate more compared to the depolarizing channel.\\


\textbf{Short range vs Long range.} We have not analyzed any model with long range interactions, nor Floquet type systems. The Quantum Kicked Top given by the following Hamiltonian

\begin{equation}
    \hat{H} = \alpha\hat{J}_{z} + \dfrac{k}{2J}\alpha\hat{J}_{x}^2\sum_{n = -\infty}^{+\infty}\delta(t-n\tau) \, ,
    \label{eq:qkt}
\end{equation}

can be realized in several ways. One, is that the model describes the dynamics of a single large spin of angular momentum $J$. The second one, is the long range interaction between $N = 2J$ single qubits, which gives the definition of the collective angular momentum operator in the Hamiltonian

\begin{equation}
   \hat{J}_{z} = \dfrac{1}{2}\sum_{i=1}^{N}\hat{\sigma}_{z}^{(i)} \, , \hat{\sigma}_{z}^{(i)} = \hat{I}_{1,i-1}\otimes\hat{\sigma}_{z}\otimes\hat{I}_{i+1,N} \, .
\end{equation}

\begin{equation}
   \hat{J}_{x}^2 = \dfrac{1}{4}\sum_{\forall i,j}^{N}\hat{\sigma}_{x}^{(i)}\hat{\sigma}_{x}^{(j)} \, , \hat{\sigma}_{x}^{(i)}\hat{\sigma}_{x}^{(j)} = \hat{I}_{1,i-1}\otimes\hat{\sigma}_{x}^{i}\otimes\hat{\sigma}_{x}^{j}\otimes\hat{I}_{j+1,N} \, .
\end{equation}

In that case, it is no longer true that the most chaotic region in the parameter space $(\alpha,k)$ is when the $\alpha \approx k$, like the models in previous section, where I assumed there could be a notion of "directionality". In fact, for any value of $\alpha$, the system is globally chaotic whenever $\alpha <<k$. In this model one can still make a bipartition and study the purity of a single spin vs the rest of the system, but the quantum channel can not be defined due to the restriction (that we have discussed also in the Bose-Hubbard model) that one can not independently evolve the basis states of the environment from the state of the system. But, the idea of the chaometer can still be applied here.\\

The QKT is an example of long-range and Floquet system.\\

The easiest model with short range interaction and Floquet is the kicked Ising model. To make a Floquet description of the Ising model, the only thing needed is to separate the unitary operator. If $\hat{U}(t) = e^{-i\hat{H}t} = e^{-i(\hat{A}+ \hat{B})t}$, then the Floquet operator after one period is just simply $\hat{F} = e^{-i\hat{A}t}e^{-i\hat{B}t}$, in the sense that the unitary propagation now reads $\ket{\psi(n)} = \hat{F}^n\ket{\psi(0)}$. If the original Hamiltonian was integrable, then the Floquet version turns out to be in most cases chaotic (due to Trotter errors), but now described by the Circular ensembles. One of the biggest differences from the time independent case, is that most of the quantities calculated as a function of the parameters, in the case of the Ising model $(g,h)$, will show periodicities and the most chaotic region in parameter space in general will no longer be when $g\approx h$.\\

\subsection{Entanglement generation = decoherence}

\textbf{Possible refined hypothesis}: With the chaometer setup in mind, any unitary process that generates an approximately random vector (pseudo-random is a more technical word, but i to define that too) will drive the purity toward its minimal value. This is equivalent to viewing the reduced dynamics as a quantum channel that closely resembles the depolarizing channel (more precisely, the reduced density matrix approximately a 2-design).\\

Due to the fact that we are interested in the dynamics on the density matrix of a single spin, maybe it is easier to define a concrete notion of directionality, as the effect of the quantum channel $\hat{\varepsilon}(t)$ on the Bloch sphere. If the Bloch sphere gets deformed in all the directions, then it is acting like the depolarizing channel with no preferred direction. With that definition and the restriction to only the purity of the chaometer in mind, one cannot distinguish chaotic Hamiltonians from completely integrable ones. Even worse, there are pseudo-chaotic evolutions\cite{Lee2024, Maziero2015RandomSampling, Gu2024Pseudochaotic, Lubkin1978Entropy} which precisely try to address the question of how to truly identify a chaotic/random unitary with limited resources, the only thing you need is that your unitary evolution gets approximately close to generate a 2-design, to fix the value of the second moment of the reduced density matrix and then obtain the value of the purity very close to 1/2. Also, it seems that simulating a 1D spin chain may not be as difficult to begin with \cite{PhysRevLett.97.157202}.\\

If, as seen by the subsystem, makes the purity decay to its lowest possible value like a depolarizing channel (we need to provide a time dependent depolarizing, it could be a RMT GOE Hamiltonian in the unitary dynamics) without any more fluctuations, then that Hamiltonian is isotropic. Is this directly connected to quantum chaos? Well not necessarily, because we are defining a property which only a subsystem can detect, we don't have the full and precise information of the global Hamiltonian, the best we can say is that it produces a reduced evolution like de depolarizing channel. Deviations from this behavior, then could be regarded as "directionality", which in the case of a Hamiltonian with quantum chaotic properties, will mean that you are departing from the most chaotic region in parameter space.\\

I was under the impression that scrambling was directly connected to the entanglement generation, but after a quick search I found that it is not the case. Scrambling is defined on the global state and local operators, but entanglement always involves correlations within the state which makes particularly the time scales very different (ADD REFERENCE, SCRAMBLING NOT ENTANGLER). So, the safest bet will be just to see, what unitary dynamics makes the full global state to become indistinguishable from a random vector, which after tracing out one part of the system, the remaining part will have an entropy of entanglement given by the Page value, therefore the purity almost 1/2.\\


\subsection{Ising results}

Even for the Ising model, there are several things which do not match quite nicely if one tries to make a comparison between total decoherence and the $\langle r\rangle$. Let's define first the total decoherence

\begin{equation}
    I(t) = \dfrac{1}{T}\int_{0}^{T}P(t)dt .
    \label{eq:decoherence}
\end{equation}

The idea behind that definition is that, for a given Hamiltonian, if it is "isotropic", then for certain parameters the decoherence will give us a number which (at the moment) seems to be minimum compared to a purity decay with higher fluctuations than this minimum. If the Hamiltonian is chaotic, we have observed that in the fully chaotic regime, there seems to exist a lowest purity decay which decays very fast. Any departure from that lowest curve, will make the value of the integral increase. In other words, if the fluctuations of the purity increase compared to this benchmark, then the value of the integral will increase. Another possibility, could be to measure the distance of the Choi state to the depolarizing channel $D(t) = 1/2 ||\hat{\rho}_{1}(t)-\hat{\sigma}||$, with $||\hat{A}||=Tr(\sqrt{{\hat{A}^\dagger}\hat{A}})$ and $\hat{\sigma}=(1/N)\hat{I}$.\\

\textbf{Idea} Can we obtain a similar curve\cite{microcanonical2024} with some property of the chaometer? Not like a quantum chaos detector, but as a more robust indicator of how chaotic in parameter space is a model which we know it exhibits chaos.

\begin{figure}
    \centering
%    \includegraphics[width=\linewidth]{figs/figs_Miguel/microcanonical.png}
    \caption{In this work, they compared the $\langle r \rangle$ with the entanglement properties of the eigenvectors. Both measures are static and independent of each other. The chaometer, combines both properties because it's a dynamical evolution which involves the eigenvalues and the entanglement properties of the eigenvectors.}
    \label{fig:placeholder}
\end{figure}

The point is to test if the model under investigation reflects RMT properties not just in the eigenvalues but at the same time with the eigenvectors. The only "possible advantage" of the chaometer is that in an "ideal experiment" one only measures properties of one spin without the need to measure more complicated observables. The second point is, what is the RMT result we will compare to? I think we can use the results from Gorin \cite{gorin_2006_dynamics, Gorinpurity2002} as a benchmark, a system connected to an RMT environment, makes the purity decay uniformly like a complete depolarizing channel. Then, any deviation from that we could say that the Hamiltonian under investigation is no longer fully chaotic.